% §II --- Related Work. Status: DRAFT. Every \cite key is in refs.bib.
\section{Related Work}\label{sec:related}

We position the paper against three lines: semantic-aware DRL resource
allocation, constrained/safe RL for wireless, and U-space separation
management. Table-free boundaries are drawn explicitly at the end of each
subsection.

\subsection{Semantic-aware DRL resource allocation}
Task-oriented and semantic communication has shifted the design objective
from bit fidelity to downstream task success~\cite{gunduz2023beyond,
xie2021deepsc}, with VQA a canonical receiver
task~\cite{xie2021mudeepsc,liu2024gosg}. Park and
Yoon~\cite{park2025transmit} transmit task-adaptive visual semantics---the
closest single-link neighbour of our evidence levels---but on one link with
no fleet contention, no airspace coupling, and no admission decision;
adaptive resource allocation for semantic-communication networks is studied
in~\cite{wang2023parameterized}, and hybrid bit/semantic transmission
reaches UAV networks in~\cite{li2026hybrid}, again without a safety
regulator in the loop. When many users and links share scarce spectrum and
compute, the allocation is naturally cast as a DRL problem. Ji
\etal~\cite{ji2024semantic} use multi-agent PPO for semantic-aware
allocation on a VQA-bearing workload, with an exhaustive upper bound and a
$100$-episode test protocol. Long \etal's
Lya-HiPPO~\cite{long2025lyahippo} couples Lyapunov-guided constraint
handling with hierarchical PPO for semantic-aware AoI minimization in
UAV-assisted networks and is the source of our flat-PPO baseline
convention. Hybrid discrete--continuous action spaces---selecting a mode
and tuning its parameters---appear in parameterized-action
RL~\cite{masson2016padrl} and its wireless
adaptations~\cite{huang2024hybrid}, and value-based methods with
conventional-optimization baselines are common in the IoT
setting~\cite{shao2024value}. Our companion
paper~\cite{zhang2026askb4tx} establishes the per-question
evidence--question complementarity that our scheduler exploits, but it
performs no multi-agent scheduling and imposes no safety constraints.
\emph{Boundary}: unlike these works we (i) route among discrete
\emph{evidence levels} rather than modulation/rate/power alone, (ii) embed
the allocation in an externally specified U-space safety scenario, and
(iii) audit which mechanism actually carries the safety load---pricing or
admission---rather than assuming the learned policy satisfies the
constraints.

\subsection{Constrained and safe RL for wireless}
Constrained MDPs~\cite{altman1999constrained} and their Lagrangian
relaxations underpin most safe-RL practice. PPO~\cite{schulman2017ppo} is
routinely extended with primal--dual updates on constraint
slack~\cite{tessler2019rcpo}, trust-region feasibility
projection~\cite{achiam2017cpo}, or Lyapunov-based
stability~\cite{chow2018lyapunov,neely2010stochastic}. In wireless, Khairy
\etal~\cite{khairy2021constrained} are the reporting-standard reference for
constrained DRL---multiplier-trajectory plots, $95\%$ confidence bands, and
an explicit zero-shot generalization section (train on $200$ devices, test
on $100/600/1000$)---and we adopt their reporting template. A recurring but
rarely-audited assumption in this literature is that the environment is
feasible, so that the dual variable converges to a bounded shadow price;
CMDP theory says exactly the opposite when it is not---under an infeasible
constraint the multiplier diverges to its
cap~\cite{altman1999constrained,chow2017risk}, and
reachability-constrained RL shows that states outside the feasible set need
a dedicated feasibility mechanism rather than a
penalty~\cite{yu2022reachability}. \emph{Boundary}: we operate the
\emph{same} dual machinery across a feasible (nominal) and an infeasible
(peak-overload) operating condition and report both pathologies
empirically---an interior shadow price in the former, a diagnostic ceiling
ratchet in the latter---and we contribute the certificate-based
escalation/admission layer as the feasibility mechanism that carries the
safety load where pricing cannot.

\subsection{U-space separation management}
U-space regulation~\cite{eu2021664} and its concept of
operations~\cite{corus2019conops,jarus2019sora,icao2005doc9854} define how
low-altitude airspace is kept safe. The SESAR JU BUBBLES
deliverable~\cite{bubbles2022d21} formalizes the separation management
service: a traffic performance envelope, a chained tactical-contact
separation minimum, a target-level-of-safety budget, and numerical
scenarios; peer-reviewed near-midair-collision modelling provides the
double-cited safety anchors~\cite{weinert2022well}. BUBBLES itself notes
that UTM has begun to adopt RL (for learning separation minima), which
legitimizes a learning approach in this domain. \emph{Boundary}: prior UTM
work treats communication performance as an exogenous requirement
(fit-for-purpose, with no hard delay numbers), whereas we close the loop---
the semantic-communication choice feeds the communication-delay term of the
separation chain, making evidence selection a first-class flight-safety
variable. To our knowledge no prior semantic-communication or DRL
resource-allocation paper grounds its entire scenario in a U-space
deliverable at the parameter level.
